\section{Finite double covers and reciprocal energy}\label{sec:geometry}

\subsection{Regular fibers and the factor of two}
Let $\Omega\subset\C$ be a bounded domain with smooth boundary, and let
$H$ be holomorphic on a neighborhood of $\overline\Omega$, nonzero on
$\partial\Omega$. When every zero in $\Omega$ is simple, set
\begin{equation}\label{eq:cover}
 X=\{(z,w):z\in\overline\Omega,\ w^2=H(z)\},\qquad
 \omega=\frac{dz}{w}.
\end{equation}
The complex structure on $X$ gives its orientation. At a zero $\alpha$,
write $H(z)=(z-\alpha)U(z)$, $U(\alpha)\ne0$. In the local coordinate
$z=\alpha+\zeta^2$,
\begin{equation}\label{eq:branchregular}
 w=\zeta\sqrt{U(\alpha+\zeta^2)},\qquad
 \omega=\frac{2\dd\zeta}{\sqrt{U(\alpha+\zeta^2)}}.
\end{equation}
Thus the branch point is regular on the cover and $\omega$ is holomorphic.

\begin{proposition}[Full energy]\label{prop:energy}
The conformally invariant energy of $\omega$ is
\begin{equation}\label{eq:energy}
 \E(H;\Omega):=\frac{\ii}{2}\int_X\omega\wedge\bar\omega
       =2\int_\Omega\frac{\dd A_z}{|H(z)|}.
\end{equation}
The factor two counts both sheets. In a local conformal coordinate,
the energy convention has no additional factor $1/2$.
\end{proposition}
\begin{proof}
Away from the finitely many branch points, $z$ is a coordinate on either
sheet and $\omega=f(z)\dd z$ with $|f|^2=|H|^{-1}$.
Since $(\ii/2)\dd z\wedge\dd\bar z=\dd A_z$, each sheet contributes
$\int_\Omega |H|^{-1}\dd A$. The branch points have area zero, and
\eqref{eq:branchregular} gives local finiteness, so the equality extends
across them.
\end{proof}

\begin{proposition}[Multiplicity and local integrability]\label{prop:multi}
If $H(z)=(z-z_0)^kU(z)$ with $U(z_0)\ne0$, its local reciprocal energy
is finite precisely when $k=1$ (or when $k=0$, meaning no zero).
A double zero gives logarithmic divergence; a zero of order $k>2$
gives a power divergence.
\end{proposition}
\begin{proof}
On a sufficiently small disk there are constants $0<u_-\le|U|\le u_+$.
The radial energy is comparable to
\begin{equation}\label{eq:radial}
 \int_0^\varepsilon r^{1-k}\dd r.
\end{equation}
This converges for $k<2$, is logarithmic for $k=2$, and diverges more
strongly for $k>2$.
The same conclusion holds on the normalization of the singular cover.
For $k=2\ell$, its two normalized branches have
$w=\pm(z-z_0)^\ell\sqrt U$, and $\omega$ has a pole of order $\ell$.
For $k=2\ell+1$, put $z-z_0=\zeta^2$; then
$\omega=2\zeta^{-2\ell}U(z_0+\zeta^2)^{-1/2}\dd\zeta$.
It is regular for $k=1$ and has a pole for each odd $k\ge3$.
These expressions also explain geometrically why the singular fiber
cannot simply be included among the smooth finite-energy fibers.
\end{proof}

For a nonzero holomorphic field on a closed bounded domain with no boundary
zero, local reciprocal integrability therefore restates simplicity of its
interior zeros. We use this observation once, as a classification of the
observable; it is not an independent simplicity estimate.

\subsection{A disk with two real branches}
In the remainder of the finite-cover analysis let
\begin{equation}\label{eq:diskpair}
 \Omega=D(C,R):=\{|z-C|<R\},\quad C\in\R,\quad
 \alpha<\beta\in\Omega\cap\R,
\end{equation}
and suppose these are exactly the two zeros of $H$, both simple. Then
\begin{equation}\label{eq:unitfactor}
 H(z)=(z-\alpha)(z-\beta)V(z),\quad
 m=\frac{\alpha+\beta}{2},\quad a=\frac{\beta-\alpha}{2}>0.
\end{equation}
The unit $V$ is holomorphic and nonvanishing on the disk and its boundary.
Choose a holomorphic square root of $V$ and put
\begin{equation}\label{eq:eta}
 y_0^2=(z-m)^2-a^2,\quad b=V^{-1/2},\quad
 \eta=\frac{dz}{y_0},\quad \omega=b(z)\eta.
\end{equation}
The two covers are identified by $w=\sqrt V\,y_0$.

\begin{lemma}[Explicit annulus coordinates]\label{lem:cylinder}
The full cover is an annulus. It is parametrized by
\begin{equation}\label{eq:Joukowski}
 z=m+a\cosh u,\qquad y_0=a\sinh u,\qquad
 u=s+\ii\theta\pmod{2\pi\ii},\qquad \eta=\dd u.
\end{equation}
Writing $\delta=m-C$, its domain is $-U(\theta)\le s\le U(\theta)$,
where
\begin{equation}\label{eq:Ushift}
 U(\theta)=\arcosh\left(
 \frac{-\delta\cos\theta+
 \sqrt{R^2+(a^2-\delta^2)\sin^2\theta}}a\right)>0.
\end{equation}
Both physical boundary components are retained.
\end{lemma}
\begin{proof}
In the nonzero coordinate $\zeta=e^u$,
$z=m+(a/2)(\zeta+\zeta^{-1})$ and
$y_0=(a/2)(\zeta-\zeta^{-1})$. The inverse is
$\zeta=(z-m+y_0)/a$, nonzero because
$(z-m+y_0)(z-m-y_0)=a^2$. Thus these are genuine coordinates on
the cover, including $u=0,\pi\ii$, where $y_0=0$ but the derivative
of the pair $(z,y_0)$ is nonzero.

The real interior roots imply $|\delta|+a<R$. Set $v=\cosh s\ge1$.
A direct calculation gives
\[
 |z-C|^2=a^2v^2+2a\delta v\cos\theta+
                \delta^2-a^2\sin^2\theta.
\]
The equation $|z-C|=R$ has the positive solution for $v$ in
\eqref{eq:Ushift}. At $s=0$, $z=m+a\cos\theta$ lies strictly inside
the disk. The quadratic expression has just one crossing for $v\ge1$,
so the domain is precisely the asserted strip. Its positive continuous
width and $2\pi$ periodicity identify an annulus, with boundaries
$s=\pm U(\theta)$. Differentiating \eqref{eq:Joukowski} gives
$dz/y_0=\dd u$, with the identity extending at the branch points.
\end{proof}

We label the boundary component over $z=C+R$ with $y_0>0$ as the plus
component. Its induced positive orientation projects counterclockwise
onto the base circle; its core period for $\eta$ is $2\pi\ii$.
The other component has the opposite sheet label, but its induced
orientation also projects counterclockwise. The sheet involution
$\iota$ satisfies $\iota^*\eta=-\eta$ and $\iota^*\omega=-\omega$.

\begin{corollary}[Full direct coordinate formula]\label{cor:direct}
If $H(\bar z)=\overline{H(z)}$, the energy can be written without a
branch singularity as
\begin{equation}\label{eq:directU}
 \E=4\int_0^\pi\int_0^1
 \frac{U(\theta)}{|V(m+a\cosh(tU(\theta)+\ii\theta))|}
          \dd t\dd\theta.
\end{equation}
\end{corollary}
\begin{proof}
In \eqref{eq:Joukowski}, the energy density is
$|b(z)|^2\dd s\dd\theta=|V(z)|^{-1}\dd s\dd\theta$.
The replacements $s\mapsto-s$ and $\theta\mapsto2\pi-\theta$
conjugate $z$, while leaving $U$ unchanged in the needed pairs.
Reality of $V$ gives the two factors of two. Finally put $s=tU(\theta)$;
the area Jacobian in $(s,\theta)$ is exactly $U(\theta)$.
\end{proof}
